\begin{frame}{자주 하는 실수: ``클러스터를 넓게 벌린다''}
  \kq{``마진 최대화'' $\neq$ ``두 클래스의 평균 위치를 멀리 벌린다''.}
  \begin{itemize}
    \item 결정 경계는 \textbf{오직} $(3,3)$과 $(-3,-3)$ 두 점의
      위치에만 의존한다.
    \item $(4,3)$을 $(100,100)$으로 옮겨도(다른 5개 점은 그대로)
      다시 학습시킨 $w, b$는 \textbf{소수점까지 완전히 동일} —
      $(100,100)$은 애초에 서포트 벡터가 아니었기 때문.
    \item ``마진 최대화''는 전체 데이터의 평균적 퍼짐이 아니라,
      \textbf{경계에 가장 가까운 소수의 점(서포트 벡터)까지의 거리}
     만을 기준으로 삼는다.
  \end{itemize}
  \vskip0.5em
  \begin{block}{실습으로 재현}
    서포트 벡터 \textbf{아닌} 점을 가장 바깥으로 크게 이동시켜
    재학습하면 $w, b$가 소수점 아래까지 동일.
    반대로 로지스틱회귀는 같은 점 이동에도 $w$가 달라진다 —
    ``모든 점이 손실에 기여''하는 모델에서는 멀리 있는 점도
    (비록 작게라도) 기울기에 영향을 주므로.
  \end{block}
\end{frame}
